The literature and descriptive evidence point to a narrower theoretical task than a full model of
genre origin or cultural diffusion. The paper needs a framework that explains why some
city-genre cells are more likely than others to cross from thin local activity into repeated local
entry. In the data, that transition is most strongly associated with local spawning flow,
target-genre active bands, and target-genre multi-band musician depth.

The empirical object in this paper is therefore local niche formation in a project-based creative
industry. Bands are projects, musicians are skilled workers, and a local scene emerges when
activity in a city-genre cell becomes thick enough to sustain repeated entry. The model is not a
theory of technological progress in the usual sense. It is a compact model of local variety
creation in which entry depends on local niche thickness, spinouts by already embedded
participants, and recombination through overlapping workers.

\subsection*{Environment}

Time is discrete. Cities are indexed by $c$, genre niches by $g$, and years by $t$. For a
city-genre cell that is still at risk of emergence, the local state at the start of period $t$ is
summarized by
\[
\left(B_{cg,t-1},\; M_{cg,t-1},\; \rho_{c,t-1}\right),
\]
where $B_{cg,t-1}$ is the number of active local bands in the focal genre, $M_{cg,t-1}$ is the
number of local multi-band musicians active in that genre, and $\rho_{c,t-1}$ is the share of
potential founders who are already embedded in the broader local scene. A band is one
differentiated local variety, and a scene emerges when the local niche becomes thick enough to
sustain repeated entry rather than isolated projects.

\subsection*{Entry}

In each period, a mass of potential founders considers launching a new band in genre $g$ in city
$c$. Founder $i$ draws embedded status $e_i \in \{0,1\}$ with
$\Pr(e_i = 1) = \rho_{c,t-1}$, an idiosyncratic fixed cost $f_i$, and a payoff shock
$\varepsilon_i$. Expected net entry value is
\begin{equation}
\Pi_{icgt}
=
\bar{\pi}_g
+ \phi_B \log(1 + B_{cg,t-1})
+ \phi_M \log(1 + M_{cg,t-1})
+ \gamma e_i
- f_i
+ \varepsilon_i,
\qquad
\phi_B,\phi_M,\gamma > 0.
\end{equation}

The term in $B_{cg,t-1}$ captures local niche thickness and labor pooling. The term in
$M_{cg,t-1}$ captures recombinant capacity through overlapping workers. The embedded-founder
premium $\gamma$ captures inherited local know-how from prior participation in the scene. Broader
city conditions are left in the background here and are absorbed empirically by fixed effects and
the paper's broader scene controls rather than being modeled explicitly.

Entry occurs when $\Pi_{icgt} \ge 0$. Aggregate entry is
\begin{equation}
N_{cgt}
=
\int \mathbf{1}\{\Pi_{icgt} \ge 0\} \, dH_{ct}(i),
\end{equation}
where $H_{ct}$ is the distribution of founder types in city $c$ at date $t$. It follows
immediately that
\[
\frac{\partial N_{cgt}}{\partial B_{cg,t-1}} > 0,
\qquad
\frac{\partial N_{cgt}}{\partial M_{cg,t-1}} > 0,
\qquad
\frac{\partial N_{cgt}}{\partial \rho_{c,t-1}} > 0.
\]

\subsection*{Dynamics and emergence}

Band stock evolves as
\begin{equation}
B_{cgt} = (1 - \delta) B_{cg,t-1} + N_{cgt},
\qquad
\delta \in (0,1).
\end{equation}

Let $N^E_{cg't}$ denote entry by embedded founders in genre $g'$. The observed spawning variable is
\begin{equation}
S_{ct} = \sum_{g'} N^E_{cg't}.
\end{equation}
Because $S_{ct}$ is indexed at the city-year level rather than the city-genre level, it is best
read as a proxy for broad local spinout intensity rather than as a purely focal-genre count.

A local scene emerges when activity crosses a viability threshold:
\begin{equation}
E_{cgt} = \mathbf{1}\{B_{cgt} \ge \bar{B}\},
\qquad
\bar{B} = 5.
\end{equation}
This threshold is operational rather than ontological. It marks the point at which a local niche
looks thick enough to sustain repeated entry, not the literal historical birth of a genre.

For cells that are still at risk, the one-period emergence hazard is
\begin{equation}
h_{cgt} = \Pr(E_{cgt} = 1 \mid E_{cg,t-1} = 0).
\end{equation}
Because $B_{cgt}$ depends on prior niche stock and current entry, this hazard is increasing in
$B_{cg,t-1}$, $M_{cg,t-1}$, and the latent embedded-founder share $\rho_{c,t-1}$.

\subsection*{Empirical mapping}

The model maps directly to the paper's three headline predictors:
\begin{align*}
B_{cg,t-1} &\leftrightarrow \text{target-genre active bands}, \\
M_{cg,t-1} &\leftrightarrow \text{target-genre multi-band musicians}, \\
\rho_{c,t-1} &\leftrightarrow \text{embedded-founder margin, proxied by } S_{c,t-1}.
\end{align*}

The reduced-form exact-year specification is therefore a compact empirical analogue of the
threshold-crossing condition:
\begin{equation}
\Pr(E_{cgt}=1 \mid E_{cg,t-1}=0)
=
G\!\left(
\alpha_{cg}
+ \lambda_t
+ \beta_S S_{c,t-1}
+ \beta_B B_{cg,t-1}
+ \beta_M M_{cg,t-1}
\right),
\end{equation}
where $G(\cdot)$ is increasing. The estimated specification adds broader scene controls and
smoothed lags, but this compact version captures the paper's core economic object.

\subsection*{Predictions and scope}

The model yields three predictions aligned with the main results. First, thicker target-genre
stock should raise later emergence because deeper local niche activity lowers the effective cost of
entry. Second, deeper overlapping-worker pools should raise later emergence because recombination
and team assembly are easier where more musicians already work across related projects. Third,
higher city-year spawning flow should predict later emergence because places that generate more
embedded spinouts also generate more viable founders than de novo entry alone would.

The interpretation remains deliberately modest. The paper does not observe knowledge transfer
directly, does not treat the emergence threshold as structural, and does not claim that spawning
is a purely genre-specific causal mechanism. The model instead rationalizes a reduced-form pattern:
new local scenes are more likely to emerge where local labor pooling is thicker, embedded
participants generate more new projects, and overlapping workers are already recombining across
related bands.
